\documentclass[11pt]{article}
\usepackage[T1]{fontenc}
\usepackage{lmodern,amsmath,amssymb,amsthm,booktabs,array}
\usepackage[hmargin=1in,vmargin=0.85in]{geometry}
\usepackage[colorlinks=true,urlcolor=blue,citecolor=blue,linkcolor=blue]{hyperref}
\newtheorem{theorem}{Theorem}
\newtheorem{lemma}[theorem]{Lemma}
\newtheorem{corollary}[theorem]{Corollary}
\newcommand{\Fp}{\mathbb F_p}
\newcommand{\HH}{\mathcal H}
\newcommand{\ZZ}{\mathbb Z}
\title{Residue-field extensions of O'Brien's examples for\protect\\
Arvind--Cameron--Ma--Maslova's commuting-graph question}
\author{Haobo Ma\thanks{ChronoAI Pte Ltd, Singapore.}\and
Wenlin Zhang\thanks{National University of Singapore; ChronoAI Pte Ltd, Singapore.}}
\date{November 15, 2026}
\begin{document}
\maketitle
\begin{abstract}
For Arvind--Cameron--Ma--Maslova's Question~4, second sentence, in
Section~9, p.~40 of arXiv:2305.07301v5 (27 July 2025)~\cite{ACMM},
we extend the examples attributed there to Eamonn O'Brien to non-isoclinic
groups of order $q^6$ with isomorphic full commuting graphs for every
prime power $q=p^f$, including characteristic two, using residue-field
coordinate charts for Heisenberg groups over two rings of order $q^2$.
Their derived groups are $(C_{p^2})^f$ and $(C_p)^{2f}$.
The common graph is chordal and perfect but neither a cograph, a split
graph nor a threshold graph; its clique
structure and degree distribution are explicit.
The case $f=1$ answers the posed all-prime question, and abelian factors
give examples of every order $p^n$ with $n\geq6$.
\end{abstract}

\section{The question and the groups}
The full commuting graph $\Gamma(G)$ of a finite group $G$ has vertex set
$G$; distinct vertices are adjacent precisely when they commute. In
particular, the identity and all other central elements are retained.
Arvind, Cameron, Ma and Maslova ask whether non-isoclinic examples with
the same commuting graph exist at every prime-sixth order~\cite[Section 9,
Question 4]{ACMM}. Their preceding discussion attributes examples for
$p=2,3,5$ to O'Brien. We give one construction for all primes.
The class-two question had already appeared in Arvind and Cameron's
\emph{Recognizing the Commuting Graph of a Finite Group}~\cite[Section~8,
p.~21]{AC}; O'Brien's examples settle it negatively, as acknowledged
in~\cite[Section~9 and p.~55]{ACMM}.
The coefficient groups for $f=1$ are familiar: Hatui, Narayanan and
Singla study their projective representations for odd $p$~\cite[Sections 1--2]{HNS}.
Our comparison uses the determinant-zero relation, not a ring isomorphism.
The same comparison works over every finite residue field and yields
the graph structure below.

For a commutative ring $T$ with identity, let $\HH(T)$ be the set $T^3$
with multiplication
\begin{equation}\label{eq:law}
 (x,y,z)(x',y',z')=(x+x',y+y',z+z'+xy').
\end{equation}
These are the upper unitriangular $3\times3$ matrices, with entries $x,y,z$
in positions $(1,2),(2,3),(1,3)$, respectively. Thus the law is associative;
the identity is $(0,0,0)$ and the inverse of $(x,y,z)$ is
$(-x,-y,-z+xy)$. Write $Z(G)$ for the centre of $G$,
$[g,h]=g^{-1}h^{-1}gh$, and $G'$ for the
subgroup generated by all such commutators. Groups $G,H$ are
\emph{isoclinic} if there are isomorphisms $G/Z(G)\to H/Z(H)$ and
$G'\to H'$ compatible with commutators of representatives. In particular,
isoclinism requires $G'\cong H'$.

Fix a prime $p$, an integer $f\geq1$, and $q=p^f$.
Choose a monic irreducible polynomial $\bar h\in\Fp[X]$ of degree $f$
and a monic lift $h\in(\ZZ/p^2\ZZ)[X]$. Set
\[
 k=\Fp[X]/(\bar h)\cong\mathbb F_q,\qquad
 R=(\ZZ/p^2\ZZ)[X]/(h),\qquad
 S=k[\varepsilon]/(\varepsilon^2).
\]
Take the coefficientwise digit section $s:k\to R$, using the unique
representatives of degree less than $f$ with coefficients in
$\{0,\ldots,p-1\}$. Monic polynomial division makes $R$ free of rank $f$
over $\ZZ/p^2\ZZ$, so reduction $\rho_R:R\to k$ has kernel $pR$
and each element of $R$ has a unique expression $s(a)+ps(b)$.
Every element of $S$ has a unique expression $a+\varepsilon b$.
Thus $|R|=|S|=q^2$, and the set bijection
\begin{equation}\label{eq:theta}
 \theta:R\longrightarrow S,\qquad s(a)+ps(b)\longmapsto a+\varepsilon b
\end{equation}
preserves zero and residues. For $f=1$, choose $h=X$; the rings reduce
to $\ZZ/p^2\ZZ$ and $\Fp[\varepsilon]/(\varepsilon^2)$.

\begin{theorem}\label{thm:main}
For every $p$ and $f$ as above, the groups $G=\HH(R)$ and $H=\HH(S)$ have order
$q^6=p^{6f}$, nilpotency class two and isomorphic full commuting graphs, but are
not isoclinic. Their centres and derived subgroups coincide, with
\[
 G'=Z(G)\cong (C_{p^2})^f,\qquad H'=Z(H)\cong (C_p)^{2f},
\]
where $C_m$ denotes the cyclic group of order $m$.
\end{theorem}

\section{A determinant-preserving bijection}
For $T=R$ or $S$, write $T^\times$ for its group of units and let
$\rho_T:T\to k$ be reduction modulo $p$ or
$\varepsilon$, respectively. Put $\nu_R(a)=ps(a)$ and
$\nu_S(a)=\varepsilon a$ for $a\in k$.
In either ring, $\nu_T$ is an additive injection onto the nonunits and
\begin{equation}\label{eq:nil}
 \nu_T(a)t=\nu_T(a\rho_T(t)),\qquad
 \nu_T(a)\nu_T(b)=0.
\end{equation}
For $R$, the equality $\ker\rho_R=pR$ and $p^2=0$ show that
$p$ times a lift depends only on its residue. This proves additivity and
the first identity; freeness over $\ZZ/p^2\ZZ$ proves injectivity.
For $S$, both identities follow from $\varepsilon^2=0$.
An element of $T$ is a unit exactly when its residue is nonzero:
in $R$, lift the residue inverse to obtain a product $1+pc$, whose inverse
is $1-pc$; in $S$, the inverse of $a+\varepsilon b$ with $a\ne0$ is
$a^{-1}-\varepsilon a^{-2}b$.

\begin{lemma}\label{lem:charts}
There is a bijection $\phi:R^2\to S^2$, with $\phi(0,0)=(0,0)$, such that
\[
 xy'-x'y=0\quad\Longleftrightarrow\quad
 \phi(x,y)_1\phi(x',y')_2-\phi(x',y')_1\phi(x,y)_2=0.
\]
Here subscripts denote the two coordinates in $S^2$.
\end{lemma}
\begin{proof}
Every vector of $T^2$ has exactly one of the following forms:
\[
\begin{array}{lll}
 \mathrm I_T(u,t)=(u,ut), &u\in T^\times,&t\in T,\\
 \mathrm {II}_T(s,u)=(u\nu_T(s),u),&s\in k,&u\in T^\times,\\
 \mathrm {III}_T(a,b)=(\nu_T(a),\nu_T(b)),&a,b\in k.&
\end{array}
\]
Indeed, a unit first coordinate gives type I, with $u=x$ and $t=u^{-1}y$.
If the first coordinate is a nonunit and the second is a unit, type II
has $u=y$, and $s$ is uniquely determined by $\nu_T(s)=u^{-1}x$.
If both coordinates are nonunits, injectivity of $\nu_T$ gives type III.
This also proves uniqueness and disjointness of the charts.

Define $\phi$ on these charts by
\begin{align*}
 \mathrm I_R(u,t)&\longmapsto \mathrm I_S(\theta(u),\theta(t)),\\
 \mathrm {II}_R(s,u)&\longmapsto \mathrm {II}_S(s,\theta(u)),\\
 \mathrm {III}_R(a,b)&\longmapsto \mathrm {III}_S(a,b).
\end{align*}
The same formulas with $\theta^{-1}$ give its inverse, and the zero
vector has type III with $a=b=0$.

For vectors $v=(x,y)$ and $w=(x',y')$, write
$D_T(v,w)=xy'-x'y$. The six unordered chart pairs give the following
complete list of zero-determinant conditions:
\begin{center}
\begin{tabular}{lll}
\toprule
First vector & Second vector & Determinant zero exactly when\\
\midrule
$\mathrm I_T(u,t)$ & $\mathrm I_T(u',t')$ & $t=t'$\\
$\mathrm I_T(u,t)$ & $\mathrm {II}_T(s,u')$ & never\\
$\mathrm I_T(u,t)$ & $\mathrm {III}_T(a,b)$ & $b=\rho_T(t)a$\\
$\mathrm {II}_T(s,u)$ & $\mathrm {II}_T(r,u')$ & $s=r$\\
$\mathrm {II}_T(s,u)$ & $\mathrm {III}_T(a,b)$ & $a=0$\\
$\mathrm {III}_T(a,b)$ & $\mathrm {III}_T(c,d)$ & always\\
\bottomrule
\end{tabular}
\end{center}
For completeness, the corresponding determinants are
\[
\begin{array}{lll}
 uu'(t'-t),&uu'(1-\nu_T(s)t),&u\nu_T(b-a\rho_T(t)),\\
 uu'\nu_T(s-r),&-u\nu_T(a),&0,
\end{array}
\]
in the order displayed. The second expression has nonzero residue;
the others give the stated conditions by cancellation of units and
injectivity of $\nu_T$. In the last two expressions we used
$\nu_T(s)\nu_T(b)=0$. Reversing a pair negates its determinant and
does not change whether it vanishes, including when $p=2$.
Equality of slopes and their residues is preserved by $\theta$.
Thus every condition in the table is preserved in both directions.
\end{proof}

\section{The graph and the isoclinism obstruction}
\begin{proof}[Proof of Theorem~\ref{thm:main}]
The order of $\HH(T)$ is $|T|^3=q^6$.
Equation~\eqref{eq:law} gives
\begin{equation}\label{eq:commutes}
 (x,y,z)(x',y',z')=(x',y',z')(x,y,z)
 \quad\Longleftrightarrow\quad xy'=x'y.
\end{equation}
If $\phi(x,y)=(X,Y)$, define
\[
 F:G\longrightarrow H,\qquad (x,y,z)\longmapsto(X,Y,\theta(z)).
\]
This is a bijection of all group elements. By Lemma~\ref{lem:charts}
and~\eqref{eq:commutes}, it preserves commutation in both directions.
Its injectivity also preserves distinctness, so $F$ is an isomorphism
of the full simple commuting graphs.

Direct substitution in the group law gives
\[
 [(x,y,z),(x',y',z')]=(0,0,xy'-x'y).
\]
All commutators lie on the central $z$-axis, and every element
$(0,0,a)$ occurs as $[(a,0,0),(0,1,0)]$. Thus
$\HH(T)'=\{(0,0,a):a\in T\}$, isomorphic to the additive group of $T$.
Conversely, commuting with $(1,0,0)$ forces $y=0$, and commuting with
$(0,1,0)$ forces $x=0$. Consequently this axis is also the whole centre.
It is nontrivial and contains all commutators, so the groups have class two.
The additive groups of $R$ and $S$ are $(C_{p^2})^f$ and
$(C_p)^{2f}$. In particular, $(0,0,1)\in G'$ has order
$p^2$, and every element of $H'$ has order dividing $p$.
The derived groups are not isomorphic, which excludes isoclinism.
All arguments apply unchanged to $p=2$.
\end{proof}

\begin{corollary}\label{cor:products}
For every prime $p$ and integer $n\geq6$, there are non-isoclinic
groups of order $p^n$ with isomorphic full commuting graphs. The two
groups in Theorem~\ref{thm:main}, and their products with any common finite
abelian group, also have isomorphic commuting graphs on their noncentral elements.
\end{corollary}
\begin{proof}
For any finite abelian group $A$, use $G\times A$ and $H\times A$.
Two elements of either product commute exactly when their first components
commute. Hence $F\times\mathrm{id}_A$ is a full commuting-graph
isomorphism. Their derived subgroups remain $G'\times\{1\}$ and
$H'\times\{1\}$ and are still nonisomorphic. Choose
$f=1$ and $A=(C_p)^{n-6}$, with the trivial group when $n=6$.
Finally, the chart map fixes the zero vector, so $F$ maps one centre onto
the other; the centres of the products are $Z(G)\times A$ and $Z(H)\times A$.
Restrict the corresponding bijections to their complements.
\end{proof}

\begin{corollary}\label{cor:structure}
The common full graph of Theorem~\ref{thm:main} is chordal and perfect,
but neither a cograph, a split graph nor a threshold graph.
Its degree distribution is
\[
\begin{array}{c|ccc}
\text{degree} &q^6-1&q^5-1&q^4-1\\
\text{multiplicity}&q^2&q^4-q^2&q^6-q^4.
\end{array}
\]
Its clique and chromatic numbers are both $q^4$. The commuting
probability is $q^{-2}+q^{-3}-q^{-5}$ and the number of conjugacy classes
is $q^4+q^3-q$.
\end{corollary}
\begin{proof}
Let $T$ be either ring. Call a vector in $T^2$ primitive if at least one
coordinate is a unit. The vertices with nonprimitive base vector form
a clique $B$ of size $q^4$, by $\nu_T(a)\nu_T(b)=0$.
The primitive vectors split into the type-I slope blocks indexed by
$t\in T$ and the type-II blocks indexed by $s\in k$. There are
$q^2+q=q(q+1)$ such blocks, each with $q^2-q$ base vectors, hence
$q^4-q^3$ group vertices after including the central coordinate.
The determinant table shows that each block is a clique, with no edges
between distinct primitive blocks. Its neighbours in $B$ have base
vectors on its one residue-field line: $b=\rho_T(t)a$ in type I,
or $a=0$ in type II. Including all central coordinates, this is a
subclique of $B$ of size $q^3$.

Every primitive vertex therefore has a clique neighbourhood. An induced
cycle of length at least four cannot contain such a vertex; nor can it
lie in $B$. Thus the graph is chordal. The same description proves
perfection directly: in any induced subgraph, first colour its clique
in $B$. For each remaining primitive block, its neighbouring vertices
in $B$ and the block together form a clique, so the block can be coloured
from a palette of size equal to that induced subgraph's clique number.
Different blocks have no edges between them. In the full graph $B$ and
each block together with its $q^3$ neighbours are cliques of size $q^4$,
and these exhaust the possibilities for a maximal clique.

Put $\pi=\nu_T(1)$. With central coordinate zero, the four base vectors
\[
 (1,0),\quad (\pi,0),\quad(0,\pi),\quad(0,1)
\]
induce a path $P_4$: the consecutive determinants vanish and the
nonconsecutive ones are $\pi$, $1$ and $\pi$, up to sign, all nonzero.
By~\cite[Proposition~10.1]{ACMM} the graph is not a cograph.
Also, two vertices from each of two distinct primitive blocks induce
$2K_2$, two disjoint edges. Such a graph has no partition into a clique
and an independent set, so our graph is not split.
Nor is it a threshold graph: if its edges were precisely the pairs
whose vertex weights sum to more than a threshold $c$, the two edges
of this $2K_2$ would give a total weight greater than $2c$, whereas
two disjoint cross nonedges would give a total at most $2c$.

A central vertex has centralizer order $q^6$. A primitive vertex has
centralizer order $q^4$, by its block and its neighbours in $B$.
For a nonzero nonprimitive vector $(\nu_T(a),\nu_T(b))$, its commuting base
vectors satisfy the single nonzero linear equation
$a\rho_T(y)-b\rho_T(x)=0$ over $k$. There are $q$ possible residue
pairs and $q^2$ lifts of each, hence centralizer order $q^5$.
Counting the three types gives the degree table.
The sum of centralizer orders is
\[
 q^2q^6+(q^4-q^2)q^5+(q^6-q^4)q^4=q^{10}+q^9-q^7.
\]
Divide by $q^{12}$ for the commuting probability. Summing
$|C_G(g)|/|G|$ over each conjugacy class gives one, so division
by $q^6$ instead gives the number of conjugacy classes.
\end{proof}

\section{Coverage and verification}
The class-two question was explicitly posed as the closing ``Conjecture''
of~\cite[Section~8, p.~21]{AC}:
\begin{quote}
Is it true that a nilpotent group of class 2 is determined up to
isoclinism by its commuting graph?
\end{quote}
It also appears at the end of Section~9 in versions~3 and~4 of~\cite{ACMM}.
Theorem~\ref{thm:main} gives an all-prime-power family answering it
negatively, extending the earlier counterexamples credited to O'Brien
in version~5.
The forward implication in~\cite[Theorem~9.1]{ACMM} is unchanged.
\begin{center}
\small
\begin{tabular}{p{0.28\linewidth}p{0.63\linewidth}}
\toprule
Source item in~\cite{ACMM} & Coverage\\
\midrule
Question 4, second sentence, \S9, p.~40 & Answered: yes for every prime
($f=1$ in Theorem~\ref{thm:main}); prior finite-prime examples:
O'Brien, reported in 2025.\\
Question 4, first sentence, same locator & Explicit family, residue-field
and abelian-factor extensions proved; general classification not covered.\\
Class-two question,~\cite[\S8, p.~21]{AC}; also~\cite[v3, p.~28;
v4, p.~38]{ACMM} & Answered: no; prior counterexamples O'Brien (2025).
Theorem~\ref{thm:main} gives examples for every $q=p^f$.\\
Question 2, \S1, p.~2 & Partial coverage proved in Corollary~\ref{cor:structure}:
this family is chordal and perfect, neither cograph, split nor threshold;
no classification.\\
Questions 1 and 3, \S1, pp.~2--3 & Not covered: general or restricted
polynomial-time recognition.\\
Question 5, \S11.2, p.~43 & Not covered: Frobenius-group classification.\\
Problem 1, \S11.3, p.~44 & Not covered: simple groups with chordal graphs.\\
\bottomrule
\end{tabular}
\end{center}

\begin{sloppypar}
The pinned Lean file
\href{https://github.com/the-omega-institute/trureturing/blob/a5ab754b3968a7da271b36fb2b1d878697d70892/D5/S3/FiniteGroups/Heisenberg/CommutingGraphQ4.lean}{\texttt{CommutingGraphQ4.lean}}
contains \texttt{graphIso} for the full graph isomorphism,
\texttt{not\_isoclinic} for the derived-subgroup obstruction, and
\texttt{result} for the existential order-$p^6$ statement for every prime,
including $2$: these concrete declarations concern $f=1$.
The full-graph and isoclinism conventions are defined in
\href{https://github.com/the-omega-institute/trureturing/blob/a5ab754b3968a7da271b36fb2b1d878697d70892/D5/S3/FiniteGroups/Heisenberg/CommutingGraph.lean}{\texttt{CommutingGraph.lean}}.
Its supporting
\href{https://github.com/the-omega-institute/trureturing/blob/a5ab754b3968a7da271b36fb2b1d878697d70892/D5/S3/FiniteGroups/Heisenberg/DeterminantCharts.lean}{\texttt{DeterminantCharts.lean}}
proves \texttt{determinantEquiv\_}\allowbreak\texttt{preserves}
for two commutative rings with \texttt{ResidueCoordinates} over a common field:
its hypotheses are a residue-respecting bijection with $k^2$, the unit
criterion, and the nilpotent multiplication identity,
and
\href{https://github.com/the-omega-institute/trureturing/blob/a5ab754b3968a7da271b36fb2b1d878697d70892/D5/S3/FiniteGroups/Heisenberg/PrimeSquareCoordinates.lean}{\texttt{PrimeSquareCoordinates.lean}}
supplies those coordinates for $\ZZ/p^2\ZZ$ and
$\Fp[\varepsilon]/(\varepsilon^2)$.
The elementary triple law and derived axis are established in
\href{https://github.com/the-omega-institute/trureturing/blob/a5ab754b3968a7da271b36fb2b1d878697d70892/D5/S3/FiniteGroups/Heisenberg/Basic.lean}{\texttt{Basic.lean}}.
The extension-field coordinate instances and the order and derived-group
calculations for $f>1$ have written proofs above.
Corollaries~\ref{cor:products} and~\ref{cor:structure} are written consequences,
not additional conclusions of \texttt{CommutingGraphQ4.result}.
The accompanying \texttt{verify.py} checks the coordinate bijection,
all chart-pair determinant conditions and derived-axis orders in exact
finite arithmetic for $p=2,3,5$. It also checks the full triple adjacency
relation for $p=2,3$. These finite checks complement the all-prime-power proof.
The additional \texttt{verify\_siblings.py} checks all ordered base-vector
pairs for $q=2,3,4,5$, the clique decomposition, degree formulas and induced
$P_4$ and $2K_2$. Its $q=4$ case uses $(\ZZ/4\ZZ)[X]/(X^2+X+1)$ and
$\mathbb F_4[\varepsilon]/(\varepsilon^2)$, with exact arithmetic.
\end{sloppypar}

{\small
\begin{thebibliography}{9}
\bibitem{AC}
V. Arvind and Peter Cameron,
\emph{Recognizing the Commuting Graph of a Finite Group},
\href{https://arxiv.org/abs/2206.01059v1}{arXiv:2206.01059v1}
(2 June 2022), Section~8, p.~21.
\bibitem{ACMM}
V. Arvind, Peter J. Cameron, Xuanlong Ma and Natalia V. Maslova,
\emph{Aspects of the commuting graph}, Journal of Algebra
\textbf{703} (2026), 4--64,
\href{https://doi.org/10.1016/j.jalgebra.2025.07.020}{doi:10.1016/j.jalgebra.2025.07.020}.
Text cited: \href{https://arxiv.org/abs/2305.07301v5}{arXiv:2305.07301v5}
(27 July 2025), Section~9, Question~4, p.~40.
\bibitem{HNS}
Sumana Hatui, E. K. Narayanan and Pooja Singla,
\emph{Projective representations of Heisenberg groups over the rings of
order $p^2$}, Journal of Group Theory \textbf{27} (2024), 413--442,
\href{https://doi.org/10.1515/jgth-2022-0033}{doi:10.1515/jgth-2022-0033}.
\href{https://arxiv.org/abs/2202.07243}{arXiv:2202.07243} (2022).
\end{thebibliography}
}
\medskip\noindent\small
\copyright\ 2026 Haobo Ma and Wenlin Zhang.
Licensed under \href{https://creativecommons.org/licenses/by/4.0/}{CC BY 4.0}.
\end{document}
